\section{Discussion}
\label{sec:discussion}
\TITLE{} uses memory as the primary admission resource because memory capacity
determines whether proof stages can safely overlap on a node. The same
framework could incorporate CPU estimates when CPU contention dominates. We
also treat partitioning and runtime scheduling separately: partitioning defines
the proof jobs, while \TITLE{} controls their execution. Finally, the
coordinator provides request-level fairness rather than explicit deadlines or
priorities, which could be added as alternative scheduling objectives.

\section{Conclusion}
\TITLE{} treats partitioned zkML proof generation as a distributed serving
problem. It separates cluster-wide job placement from node-local stage
admission, uses historical resource profiles to control memory safely, and
reuses expensive setup artifacts across requests. Our evaluation shows that
partitioning exposes heterogeneous proof jobs with strongly time-varying memory
demand and that those jobs scale effectively across prover nodes.
\textbf{[Insert final concurrent-serving headline result.]}
